\documentclass[12pt,letterpaper]{article}
\usepackage[margin=0.5in]{geometry}
\usepackage{array}
\usepackage{amsmath}
\usepackage{graphicx}
\usepackage{xcolor}
\usepackage{tikz}
\renewcommand{\familydefault}{\sfdefault}
\setlength{\parindent}{0pt}
\setlength{\parskip}{0.3em}
\setlength{\fboxsep}{3pt}

\newcolumntype{C}[1]{>{\centering\arraybackslash}m{#1}}
\newcommand{\onetape}[2]{%
  \begin{tikzpicture}[x=0.10in,y=0.11in,baseline=-0.5ex]
    \draw (0,0) rectangle (#1,1);
    \foreach \x in {1,...,#1} {\draw (\x,0) -- (\x,1);}
    \fill[blue!14] (0,0) rectangle (#2,1);
    \draw (0,0) rectangle (#1,1);
    \foreach \x in {1,...,#1} {\draw (\x,0) -- (\x,1);}
  \end{tikzpicture}%
}
\newcommand{\tapediagrams}{\onetape{8}{2}\hspace{0.35em}\onetape{4}{1}}
\newcommand{\fractionnumberline}{%
  \begin{center}
    \begin{tikzpicture}[x=1in]
      \draw[<->,line width=0.45pt] (0,0) -- (6.9,0);
    \end{tikzpicture}
  \end{center}
}
\newcommand{\answerbox}[1]{%
  \begin{center}
    {\bfseries #1}\\
    \vspace*{\fill}
    {\Large Write your answer\\ in your Class Notes}
  \end{center}
  \vspace*{\fill}
}

\begin{document}

\begin{center}
    {\Large\bfseries Assignment 4: Fractions}\\[4pt]
    {\large\bfseries DO NOT WRITE ON THIS PAPER}\\[4pt]
\end{center}
\vspace{-0.8em}

\noindent\hfill\fbox{%
  \begin{minipage}[t][0.34\textheight][t]{.95\textwidth}
    \begin{center}
      \textbf{Part A}
    \end{center}

    On your whiteboard, draw the table below.

    Flip two cards. The first card is your numerator and the second card is your denominator. Find an
    equivalent fraction to your original fraction and draw tape diagrams to represent both fractions.

    \begin{center}
    \renewcommand{\arraystretch}{1.55}
    \begin{tabular}{|C{0.10\textwidth}|C{0.10\textwidth}|C{0.16\textwidth}|C{0.23\textwidth}|C{0.21\textwidth}|}
      \hline
      \rule{0pt}{2.7em}\textbf{Card 1} & \textbf{Card 2} & \textbf{Fraction: }$\frac{\text{card1}}{\text{card2}}$ &
      \textbf{Equivalent Fraction} & \textbf{Tape Diagrams}\\
      \hline
      \rule{0pt}{2.8em}2 & 8 & $\frac{2}{8}$ & $\frac{2\div2}{8\div2}=\frac{1}{4}$ & \tapediagrams
      \\
      \hline
      \rule{0pt}{1.9em} & & & & \\
      \hline
    \end{tabular}
    \end{center}

    On your paper write one example and your answer to: What does it mean for two fractions to be
    equivalent?

    \vfill
  \end{minipage}%
}\hfill

\noindent\hfill\fbox{%
  \begin{minipage}[t][0.48\textheight][t]{.95\textwidth}
    \begin{center}
      \textbf{Part B}
    \end{center}

    Use the process from part A to create two fractions. Convert each fraction so they have common denom-
    inator.

    Add to evaluate the sum Fraction 1 + Fraction 2. Subtract to evaluate the difference Fraction 1 -- Fraction
    2.

    \begin{center}
    \renewcommand{\arraystretch}{1.65}
    \setlength{\tabcolsep}{3pt}
    \scriptsize
    \begin{tabular}{|C{0.08\textwidth}|C{0.08\textwidth}|C{0.12\textwidth}|C{0.14\textwidth}|C{0.14\textwidth}|C{0.15\textwidth}|C{0.16\textwidth}|}
      \hline
      \textbf{Fraction 1} & \textbf{Fraction 2} & \textbf{Common Denom.} &
      \textbf{F1 w/ Common Denom.} & \textbf{F2 w/ Common Denom.} &
      \textbf{Sum (F1 + F2)} & \textbf{Difference (F1 - F2)}\\
      \hline
      $\dfrac{2}{3}$ & $\dfrac{5}{2}$ & 6 &
      $\dfrac{2}{3}\cdot\dfrac{2}{2}=\dfrac{4}{6}$ &
      $\dfrac{5}{2}\cdot\dfrac{3}{3}=\dfrac{15}{6}$ &
      $\dfrac{4}{6}+\dfrac{15}{6}=\dfrac{19}{6}$ &
      $\dfrac{4}{6}-\dfrac{15}{6}=-\dfrac{11}{6}$
      \\
      \hline
      \rule{0pt}{2.2em} & & & & & & \\
      \hline
    \end{tabular}
    \normalsize
    \end{center}

    \vfill
  \end{minipage}%
}\hfill

\newpage

\noindent\hfill\fbox{%
  \begin{minipage}[t][0.40\textheight][t]{.95\textwidth}
    \begin{center}
      \textbf{Part C}
    \end{center}

    Place the following numbers on a number line in the correct order (from smallest to biggest) and with
    appropriate spacing:
    \[
      0,\;1,\;\frac{1}{4},\;\frac{1}{3},\;\frac{2}{3},\;\frac{1}{2},\;\frac{2}{5},\;\frac{3}{4}
    \]

    \vspace{1.5em}
    \fractionnumberline
    \vspace{1.5em}

    \emph{Hint: Think about equivalent fractions}
    \vfill
  \end{minipage}%
}\hfill

\noindent\hfill\fbox{%
  \begin{minipage}[t][0.46\textheight][t]{.95\textwidth}
    \begin{center}
      \textbf{Question 1}
    \end{center}

    \textbf{Big Question:} Draw a tape diagram to represent $11/6$. Lets say the unit represents a 6-pack of soda.

    \begin{enumerate}
      \item[a)] Describe what the denominator tells us about the pack of soda and how you drew that in your
      diagram.
      \item[b)] Describe what the numerator tells us about packs of soda and how you drew that in your diagram.
      \item[c)] The numerator is bigger than the denominator for this fraction. What does this mean in terms of
      packs of soda, and how did you draw this in your diagram?
    \end{enumerate}

    Use the following mathematical terms in your descriptions: denominator, numerator, unit
    \vfill
  \end{minipage}%
}\hfill

\newpage

\noindent\hfill\fbox{%
  \begin{minipage}[t][0.42\textheight][t]{.70\textwidth}
    \begin{center}
      \textbf{Question 2 Work}
    \end{center}

    You walk $\dfrac{1}{8}$ miles to a shoe store. Then you walk another $\dfrac{1}{3}$ miles to a clothing
    store. How many miles have you walked in all?

    Your answer should include a fraction in simplest form and be written in a complete sentence.
    \vfill
  \end{minipage}%
}%
\fbox{%
  \begin{minipage}[t][0.42\textheight][t]{0.24\textwidth}
    \answerbox{Question 2 Answer}
  \end{minipage}%
}\hfill

\noindent\hfill\fbox{%
  \begin{minipage}[t][0.42\textheight][t]{.70\textwidth}
    \begin{center}
      \textbf{Question 3 Work}
    \end{center}

    Maya finished her English assignment in $\dfrac{2}{5}$ hours. Then she completed her
    math assignment in $\dfrac{1}{2}$ hours. How much more time did Maya spend on her
    math assignment?

    Your answer should include a fraction in simplest form and be written in a complete sentence.
    \vfill
  \end{minipage}%
}%
\fbox{%
  \begin{minipage}[t][0.42\textheight][t]{0.24\textwidth}
    \answerbox{Question 3 Answer}
  \end{minipage}%
}\hfill

\newpage

\begin{center}
    \textbf{Additional Space}
\end{center}

\vspace{0.4em}

\noindent\makebox[\textwidth][c]{%
\fbox{%
  \rule{0pt}{0.20\textheight}\rule{7.35in}{0pt}%
}%
}

\vspace{0.45em}

\noindent\makebox[\textwidth][c]{%
\fbox{%
  \rule{0pt}{0.20\textheight}\rule{7.35in}{0pt}%
}%
}

\vspace{0.45em}

\noindent\makebox[\textwidth][c]{%
\fbox{%
  \rule{0pt}{0.20\textheight}\rule{7.35in}{0pt}%
}%
}

\vspace{0.45em}

\noindent\makebox[\textwidth][c]{%
\fbox{%
  \rule{0pt}{0.20\textheight}\rule{7.35in}{0pt}%
}%
}

\end{document}
